% TOP_PROBLEM: {"id": 1399, "title": "List Coloring Conjecture", "category_id": 3, "category": {"id": 3, "name": "graph_theory", "display_name": "Graph Theory", "description": "Problems involving graphs, networks, and their properties.", "slug": "graph-theory", "order_index": 3, "created_at": "2026-07-31T15:26:25.671Z"}, "status": "open"}
% FIRST_POSTED: null
% ENTRY_KIND: statement_only
% Imported from: batch07.tex; UnsolvedMath ID 1399
\documentclass[11pt]{article}
\usepackage[margin=25mm]{geometry}
\usepackage{fontspec}
\setmainfont{Latin Modern Roman}
\usepackage{amsmath,amssymb,mathtools}
\usepackage{unicode-math}
\setmathfont{Latin Modern Math}
\usepackage{xurl,hyperref}
\hypersetup{colorlinks=true,urlcolor=blue}
\setcounter{secnumdepth}{0}
\setlength{\parindent}{0pt}
\setlength{\parskip}{5pt}
\setlength{\emergencystretch}{3em}
\newcommand{\sref}[2]{\hyperlink{source:#1:#2}{[S#2]}}
\newcommand{\cataloguescope}[1]{\paragraph{Recorded target.}#1}
\begin{document}
% UnsolvedMath ID 1399; source catalogue code GRAPH-012
\setcounter{section}{1398}
\section{List Coloring Conjecture}\label{1399}
{\small\sffamily UnsolvedMath ID 1399 · Graph Theory}
\subsection{Definitions and mathematical statement}

\paragraph{Shared notation.}$\mathbb N=\{1,2,\ldots\}$, and $\mathbb Z,\mathbb Q,\mathbb R,\mathbb C$ denote the integers, rationals, reals and complex numbers. Any other conventions are specified locally.


Let $G=(V,E)$ be a finite undirected loopless multigraph; parallel edges are allowed. Two different edges are adjacent when they share an endpoint. A proper edge coloring assigns colors to edges so that adjacent edges have different colors. The chromatic index $\chi'(G)$ is the smallest number of colors allowing such a coloring. For a choice of finite color lists $L(e)$, a proper list edge coloring must additionally assign to each edge a color from its own list. The list chromatic index $\chi'_{\ell}(G)$ is the least nonnegative integer $k$ such that every assignment of lists with $|L(e)|\ge k$ has a proper list edge coloring. Both indices are zero when $E$ is empty.
\paragraph{Statement recorded in the supplied source.}
For every finite loopless multigraph $G$, is
\[
 \chi'_{\ell}(G)=\chi'(G)?
\]
Equivalently, is a proper edge coloring always possible when each edge is given an arbitrary list of $\chi'(G)$ colors? The statement concerns every graph, including non-bipartite graphs. \sref{1399}{2}
\subsection{Short English statement}
Does allowing a different list of permitted colors on each edge ever require more colors than ordinary proper edge coloring?
\subsection{Sources}
\begin{description}
\item[\hypertarget{source:1399:1}{S1}] ULAM.AI and the UnsolvedMath Contributors, UnsolvedMath: A Curated Collection of Open Mathematics Problems, record 1399 (GRAPH-012). CC BY 4.0. \url{https://huggingface.co/datasets/ulamai/UnsolvedMath}.
\item[\hypertarget{source:1399:2}{S2}] Fotis Iliopoulos and Alistair Sinclair, Efficiently List-Edge Coloring Multigraphs Asymptotically Optimally, arXiv:1812.10309v3 (2021), Section 1, list chromatic index and List Coloring Conjecture preceding Theorem 1.2. \url{https://arxiv.org/pdf/1812.10309v3}.
\end{description}
\label{end:1399}
\end{document}
